% !TeX root = main.tex
\documentclass[12pt, a4paper]{article}

\usepackage[utf8]{inputenc}
\usepackage[T1]{fontenc}
\usepackage[margin=2.5cm]{geometry}
\usepackage{amsmath, amssymb, amsfonts}
\usepackage{graphicx}
\usepackage{xcolor}
\usepackage{booktabs}
\usepackage{tabularx}
\usepackage{longtable}
\usepackage{multirow}
\usepackage{array}
\usepackage{hyperref}
\usepackage{cite}
\usepackage{url}
\usepackage{microtype}
\usepackage{enumitem}
\usepackage{fancyhdr}
\usepackage{titlesec}
\usepackage{caption}
\usepackage{subcaption}
\usepackage{float}
\usepackage{parskip}
\usepackage{tcolorbox}
\usepackage{mdframed}
\usepackage{tikz}
\usetikzlibrary{fit, positioning, arrows.meta, backgrounds, shadows.blur, decorations.pathreplacing}

\graphicspath{{./}{figures/}}
\providecommand{\doi}[1]{\url{https://doi.org/#1}}

\definecolor{AccentBlue}{HTML}{2563EB}
\definecolor{AccentGreen}{HTML}{059669}
\definecolor{AccentAmber}{HTML}{D97706}
\definecolor{AccentRed}{HTML}{DC2626}
\definecolor{AccentViolet}{HTML}{7C3AED}
\definecolor{CampusTeal}{HTML}{0D9488}
\definecolor{PanelGray}{HTML}{F8FAFC}
\definecolor{BorderGray}{HTML}{E2E8F0}
\definecolor{TextDark}{HTML}{1E293B}
\definecolor{TextMuted}{HTML}{64748B}

\definecolor{GridOrange}{HTML}{F97316}
\definecolor{GridAmber}{HTML}{FBBF24}
\definecolor{SensorBlue}{HTML}{3B82F6}
\definecolor{SensorLightBlue}{HTML}{BFDBFE}
\definecolor{AIGreen}{HTML}{10B981}
\definecolor{AILightGreen}{HTML}{A7F3D0}
\definecolor{GovernRed}{HTML}{EF4444}
\definecolor{GovernLightRed}{HTML}{FECACA}
\definecolor{NoosphereViolet}{HTML}{8B5CF6}
\definecolor{NoosphereLightViolet}{HTML}{DDD6FE}
\definecolor{NodeBg}{HTML}{F8FAFC}
\definecolor{LayerBg}{HTML}{F1F5F9}

\hypersetup{
  colorlinks=true, linkcolor=AccentBlue,
  citecolor=AccentBlue, urlcolor=AccentBlue,
  pdftitle={EnergyMind Phase 1 Report},
  pdfauthor={Anonymous Team — TSYP14 IEEE}
}

\pagestyle{fancy}
\fancyhf{}
\fancyhead[L]{\small\color{TextDark!70}\textbf{EnergyMind} — Campus Microgrid AI Trust Layer}
\fancyhead[R]{\small\color{TextDark!70}TSYP14 IEEE PES \& IES — Phase 1}
\fancyfoot[C]{\small\color{TextDark!50}\thepage}
\renewcommand{\headrulewidth}{0.4pt}

\titleformat{\section}{\large\bfseries\color{AccentBlue}}{}{0em}{}[\color{AccentBlue}\hrule\vspace{2pt}]
\titleformat{\subsection}{\normalsize\bfseries\color{TextDark}}{\thesubsection}{0.8em}{}
\titleformat{\subsubsection}{\normalsize\bfseries\color{AccentGreen}}{}{0em}{}

\captionsetup{font=small, labelfont=bf, labelsep=period}

\tcbset{
  resultbox/.style={
    colback=AccentGreen!5, colframe=AccentGreen!60,
    boxrule=0.8pt, arc=4pt, left=8pt, right=8pt, top=6pt, bottom=6pt
  },
  warnbox/.style={
    colback=AccentAmber!5, colframe=AccentAmber!60,
    boxrule=0.8pt, arc=4pt, left=8pt, right=8pt, top=6pt, bottom=6pt
  }
}

\begin{document}

% ─────────────────────────────────────────────────────────────────────────────
%  TITLE PAGE
% ─────────────────────────────────────────────────────────────────────────────

\begin{titlepage}
\centering
\vspace*{2cm}

{\Huge\bfseries\color{AccentBlue} EnergyMind}\\[0.4cm]
{\Large\bfseries An AI Trust Layer for Campus Microgrid\\Cyber-Physical Anomaly Detection}\\[1.0cm]

\begin{tcolorbox}[colback=AccentBlue!8, colframe=AccentBlue!40, boxrule=0.8pt, arc=6pt,
  width=0.85\linewidth]
\centering
{\large\bfseries TSYP14 Technical Challenge — IEEE PES \& IES Chapters}\\[0.3cm]
{\normalsize Phase 1 Report: Ideation, Data Exploration \& Preliminary Design}\\[0.2cm]
{\small\color{TextDark!60} Tunisia Section · IEEE TSYP14 Congress 2026}
\end{tcolorbox}

\vspace{1.5cm}

\begin{tabular}{ll}
\toprule
\textbf{Challenge Track} & IEEE PES (Power \& Energy) + IES (Industrial Electronics)\\
\textbf{Submission Phase} & Phase 1 — Ideation \& Data Exploration\\
\textbf{System Under Study} & University Campus Microgrid (20 kV/0.4 kV, 7-bus)\\
\textbf{Repository} & \url{https://github.com/anonymous-tsyp14/energymind}\\
\textbf{Version} & EnergyMind V2 · 5 Seeds · 21/21 Tests Passing\\
\bottomrule
\end{tabular}

\vspace{2cm}

\begin{tcolorbox}[resultbox, width=0.85\linewidth]
\centering\bfseries Key Phase 1 Results\\[0.3cm]
\normalfont\small
\begin{tabular}{lr|lr}
Event Recall & \textbf{58.1\%} (CI95: 54.5–61.7\%) &
Precision & \textbf{91.2\%}\\
Calibrated ECE & \textbf{0.038} (target $\le$0.05 ✓) &
Surrogate MAE & \textbf{0.096 kW} ($<$0.1 kW ✓)\\
Diagnostic Acc. (decided) & \textbf{100\%} (5 seeds) &
False Alarms / Day & \textbf{0.173}\\
\end{tabular}
\end{tcolorbox}

\vfill
{\small\color{TextDark!50} Anonymous submission — no institutional affiliation disclosed per double-blind policy}
\end{titlepage}

% ─────────────────────────────────────────────────────────────────────────────
%  TABLE OF CONTENTS
% ─────────────────────────────────────────────────────────────────────────────
\tableofcontents
\newpage

% ─────────────────────────────────────────────────────────────────────────────
\section{Executive Summary}
% ─────────────────────────────────────────────────────────────────────────────

\textbf{EnergyMind} is a hybrid AI trust layer that stands between raw microgrid telemetry and autonomous control decisions. It addresses a genuine engineering gap: as microgrids integrate more IoT sensors, distributed generation, and remote automation, the risk of incorrect AI decisions—driven by uncalibrated probability estimates or adversarial sensor data—grows substantially. EnergyMind solves this through a principled multi-layer architecture:

\begin{enumerate}[leftmargin=*, label=\textbf{\arabic*.}]
  \item \textbf{Physics-grounded detection} — five normalized conservation residuals catch both physical anomalies and data falsification
  \item \textbf{Formal probability calibration} — isotonic regression ensures that a 90\% confidence claim is empirically 90\% correct (ECE = 0.038)
  \item \textbf{Honest uncertainty} — the system abstains on 50.6\% of ambiguous alerts rather than forcing a wrong answer; 100\% accuracy on decided cases
  \item \textbf{Privacy-preserving collective defense} — Noosphere shares anomaly signatures across a peer-to-peer network of energy nodes with $\varepsilon$-differential privacy and $k$-anonymity, not just with a single remote partner
\end{enumerate}

All results are reproducible: fixed seeds, independent test/calibration splits, 21/21 automated tests passing.

% ─────────────────────────────────────────────────────────────────────────────
\section{Scope Definition}
% ─────────────────────────────────────────────────────────────────────────────

\subsection{System Monitored}

EnergyMind monitors a \textbf{university campus microgrid} operating at 20\,kV/0.4\,kV. The system boundary covers all generation, storage, load, and interconnection equipment, sampled at 15-minute resolution.

\begin{table}[H]
\centering
\caption{Microgrid Electrical Parameters}
\begin{tabularx}{\linewidth}{llX}
\toprule
\textbf{Component} & \textbf{Parameter} & \textbf{Value / Specification}\\
\midrule
Grid Interconnection & PCC voltage & 20.0\,kV (medium voltage, STEG grid)\\
Transformer & Rating / ratio & 250\,kVA, 20\,kV/0.4\,kV, Dyn11, $u_k = 4\%$\\
LV Busbar & Nominal voltage & 400\,V (line-to-line, 7-bus network)\\
PV Array & Peak capacity & 50\,kWp rooftop installation\\
BESS & Energy / power & 100\,kWh / 25\,kW lithium-ion\\
Load feeders & Configuration & 3 branches: Admin, Labs\,\&\,HVAC, Residence\,\&\,EV\\
Telemetry & Resolution & 15-minute intervals, real-time SCADA monitoring\\
\bottomrule
\end{tabularx}
\end{table}

Figure~\ref{fig:single_line} shows the full single-line diagram generated by pandapower.

\begin{figure}[H]
  \centering
  \includegraphics[width=0.90\linewidth]{figures/fig01_single_line_diagram.png}
  \caption{Campus microgrid 7-bus single-line diagram. Nodes represent grid PCC (gray), transformer (violet), LV busbar (blue), load feeders (cyan), PV array (amber), and BESS (green). Red triangles indicate voltage measurement points $V_1$ (PCC) and $V_2$ (inverter), used by the physics residual engine.}
  \label{fig:single_line}
\end{figure}

\subsection{Eight-Class Anomaly Taxonomy}

A critical design principle is the strict separation between \emph{physical ground truth} and \emph{sensor observations}. Under a cyber-falsification attack, the physical grid may be normal while sensors report falsified measurements. EnergyMind tracks both layers simultaneously.

\begin{table}[H]
\centering
\caption{Eight-Class Cyber-Physical Anomaly Taxonomy}
\begin{tabular}{clll}
\toprule
\textbf{Code} & \textbf{Class} & \textbf{Layer} & \textbf{Detection Challenge}\\
\midrule
C0 & \texttt{normal} & Baseline & Baseline operating condition\\
P2 & \texttt{pv\_fault} & Physical & PV string degradation, partial shading\\
G1 & \texttt{grid\_event} & Physical & Voltage sag/swell at PCC\\
L1 & \texttt{load\_peak} & Physical & Unexpected high-load transient\\
M1/M2 & \texttt{pv\_sensor\_fault} & Sensor & PV meter drift or calibration bias\\
V1 & \texttt{voltage\_sensor\_fault} & Sensor & PCC voltage transducer offset\\
B1 & \texttt{bess\_sensor\_fault} & Sensor & Battery SoC coulomb counter drift\\
A1 & \texttt{coordinated\_falsif.} & Cyber & Multi-sensor stealth falsification\\
\bottomrule
\end{tabular}
\end{table}

\begin{tcolorbox}[warnbox]
\small\textbf{Key integrity property:} Classes P2 (\texttt{pv\_fault}) and A1 (\texttt{coordinated\_falsification}) produce \emph{identical sensor observations} but different physical ground truths. Separating them requires querying the independent inverter sensor — this drives the Active Evidence Acquisition mechanism.
\end{tcolorbox}

% ─────────────────────────────────────────────────────────────────────────────
\section{Objectives}
% ─────────────────────────────────────────────────────────────────────────────

EnergyMind addresses all four TSYP14 Challenge Goals:

\begin{table}[H]
\centering
\caption{Challenge Goals and EnergyMind's Response}
\begin{tabularx}{\linewidth}{clX}
\toprule
\textbf{Goal} & \textbf{Challenge Requirement} & \textbf{EnergyMind Implementation}\\
\midrule
1 & Evidence Fusion \& Fault Detection & Physics residuals + IsoForest + CUSUM hybrid fusion\\
2 & Trust \& Uncertainty Quantification & Isotonic calibration (ECE=0.038), selective abstention, Active Evidence Acquisition\\
3 & Governance \& Explainability & SHA-256 audit chain, SHAP per-alert attribution, expected-loss policy\\
4 & Collective Defense & Noosphere: peer-to-peer ε-DP knowledge transfer across multiple energy nodes\\
\bottomrule
\end{tabularx}
\end{table}

\textbf{Core engineering commitments:}
\begin{itemize}
  \item All announced functions exist and pass automated unit tests (21/21)
  \item All metrics computed from held-out test sets with deterministic seeds
  \item Every model choice is justified empirically against alternatives (see \texttt{docs/MODEL\_SELECTION.md})
  \item Limitations are documented honestly — no metric inflation
\end{itemize}

% ─────────────────────────────────────────────────────────────────────────────
\section{Data \& Anomaly Simulation Plan}
% ─────────────────────────────────────────────────────────────────────────────

\subsection{Real Baseline Data}

Real grid fault/anomaly-labeled data is rare. EnergyMind uses two real campus datasets as its foundation:

\begin{table}[H]
\centering
\caption{Real Baseline Datasets Used}
\begin{tabularx}{\linewidth}{lXlX}
\toprule
\textbf{Dataset} & \textbf{Description} & \textbf{Volume} & \textbf{Role}\\
\midrule
Campus PV measurements & 5-min output from campus rooftop array & 238,464 rows & Solar baseline normalization\\
Campus electricity invoices & Monthly bills 2022--2024 & 30 records & Anchor load consumption to real scale\\
\bottomrule
\end{tabularx}
\end{table}

\subsection{Cross-Validation with International References}

\begin{table}[H]
\centering
\caption{External Reference Datasets for Validation}
\begin{tabularx}{\linewidth}{llX}
\toprule
\textbf{Source} & \textbf{Organization} & \textbf{Validation Result}\\
\midrule
PVGIS API \cite{pvgis} & EU Commission JRC & 69,752.9\,kWh/yr for 50\,kWp at 36.8°N 10.2°E (matches campus invoices ±8.1\%)\\
NASA POWER \cite{nasa_power} & NASA Langley & Hourly GHI + $T_{2m}$ reanalysis (MERRA-2) for weather normalization\\
SimBench 1-LV-semiurb4 \cite{simbench} & Fraunhofer IEE & Line impedances $R/X$ for the 7-bus network\\
DOE Commercial Buildings \cite{doe_bldg} & US Dept. Energy & Office, Lab, HVAC load archetypes\\
UAH ICS Cyber Dataset \cite{uah_ics} & Tommy Morris et al. & Cyber attack taxonomy for A1 class injection\\
\bottomrule
\end{tabularx}
\end{table}

Figure~\ref{fig:powerflow} validates the pandapower AC power flow model against physical constraints.

\begin{figure}[H]
  \centering
  \includegraphics[width=0.95\linewidth]{figures/fig02_powerflow_states.png}
  \caption{Pandapower AC power flow analysis for the 7-bus campus network. \textbf{Left:} PCC busbar voltage as a function of PV output under three load regimes (30, 45, 60\,kW). Dashed lines show IEEE 1547 voltage limits (360–440\,V). \textbf{Right:} Net grid import vs.\ PV generation; the zero-crossing marks the self-sufficiency operating point where the campus fully covers its load from local resources.}
  \label{fig:powerflow}
\end{figure}

\subsection{Anomaly Injection Methodology}

Anomalies are injected by \texttt{faults\_v2.py} using deterministic multi-seed reproducibility. For each class $c_k$, non-overlapping time windows are selected and both physical consequences and sensor corruptions are computed independently:

\begin{equation}
\mathbf{y}_{\text{obs}}(t) = \mathcal{T}\!\left(\mathbf{x}_{\text{true}}(t),\,\boldsymbol{\theta}_{\text{fault}}(t)\right) + \boldsymbol{\epsilon}(t),
\quad \boldsymbol{\epsilon} \sim \mathcal{N}(\mathbf{0},\,\boldsymbol{\Sigma}_{\text{noise}})
\end{equation}

where $\mathcal{T}$ is the corruption operator. Under normal conditions $\mathcal{T}=\mathbf{I}$; under cyber-attacks it falsifies measurements while $\mathbf{x}_{\text{true}}$ remains normal.

All injected events are logged to \texttt{outputs/calibration\_event\_registry\_v2.csv} and \texttt{outputs/test\_event\_registry\_v2.csv} with timestamps, class labels, and injection parameters.

% ─────────────────────────────────────────────────────────────────────────────
\section{System Architecture \& Preliminary Design}
% ─────────────────────────────────────────────────────────────────────────────

\subsection{Global Architecture Overview}

Figure~\ref{fig:architecture} presents the complete EnergyMind five-layer architecture.

\begin{figure}[H]
  \centering
  \resizebox{0.98\linewidth}{!}{\input{architecture.tex}}
  \caption{EnergyMind global system architecture. Five distinct layers: (1) Physical 7-bus microgrid with PV, BESS, and campus feeders; (2) Sensor and telemetry acquisition; (3) AI Trust Layer with hybrid detection, Bayesian diagnosis, and isotonic calibration; (4) Governance with SHA-256 audit chain and SHAP explainability; (5) Noosphere federated multi-node collective defense network. Color coding: amber = power flow, blue = sensor telemetry, green = AI processing, red = governance, violet = federated network.}
  \label{fig:architecture}
\end{figure}

\subsection{Pillar 1 — Physics Residuals ($Z \in \mathbb{R}^5$)}

EnergyMind computes five physics-grounded normalized residuals from raw telemetry:

\begin{align}
r_{\text{bal}}(t) &= \frac{P_{\text{pv}} + P_{\text{grid}} + P_{\text{bess}} - P_{\text{load}}}{\sigma_{\text{bal}}} & &\text{(power conservation)}\\
r_{\text{pvexp}}(t) &= \frac{P_{\text{pv}} - \hat{P}_{\text{pv}}^{\text{LGB}}(G, T)}{\sigma_{\text{pvexp}}} & &\text{(PV model residual, LightGBM)}\\
r_{\text{vpcc}}(t) &= \frac{V_{\text{pcc}} - \hat{V}_{\text{pcc}}^{\text{surrogate}}}{\sigma_{\text{vpcc}}} & &\text{(bus voltage deviation)}\\
r_{\text{inv}}(t) &= \frac{P_{\text{inv}} - P_{\text{dc}}\,\eta_{\text{inv}}}{\sigma_{\text{inv}}} & &\text{(inverter consistency, queried on demand)}\\
r_{\text{vinv}}(t) &= \frac{V_{\text{inv}} - V_{\text{pcc}} - \Delta V_{\text{drop}}}{\sigma_{\text{vinv}}} & &\text{(inverter voltage drop)}
\end{align}

\subsection{Pillar 2 — Pandapower Surrogate (Speed vs.\ Accuracy)}

Running full Newton-Raphson AC power flow at every telemetry step is computationally expensive ($\approx$200\,ms/call). A linear surrogate $\hat{V}_{\text{pcc}} = \beta_0 + \beta_1 P_{\text{net}} + \beta_2 Q_{\text{net}}$ is trained on 2,000 pandapower evaluations and validated on 500 independent states.

\begin{figure}[H]
  \centering
  \includegraphics[width=0.95\linewidth]{figures/fig03_surrogate_validation.png}
  \caption{Pandapower linear surrogate validation on 500 independent power flow states. \textbf{Left:} Scatter plot of surrogate predictions vs.\ full Newton-Raphson ground truth ($R^2 = 1.000$, MAE = 0.096\,kW). The identity line is shown in dashed red; all points cluster tightly around it. \textbf{Right:} Error quantile distribution (P50 = 0.048\,kW, P95 = 0.249\,kW, P99 = 0.421\,kW). The surrogate accelerates inference by 150× while remaining within 0.5\,kW of the full AC solution.}
  \label{fig:surrogate}
\end{figure}

\subsection{Pillar 3 — Hybrid Detection Fusion Gate}

The alert trigger $\mathcal{A}(t) \in \{0,1\}$ is:

\begin{equation}
\mathcal{A}(t) = \underbrace{\mathbb{I}\!\left(\|\mathbf{Z}\|_2 > \gamma_{\text{phys}}\right)}_{\text{physics plausibility}} \wedge \left[\underbrace{\mathbb{I}_{\text{IsoForest}}(\mathbf{Z})}_{\text{spatial outlier}} \vee \underbrace{\mathbb{I}_{\text{CUSUM}}(\mathbf{Z})}_{\text{change-point}}\right]
\end{equation}

This design enforces that alarms cannot be raised without a physical residual violation, while allowing either the unsupervised spatial outlier model or the sequential change-point detector to provide the trigger.

\subsection{Pillar 4 — Isotonic Calibration (Goal 2)}

Raw Gaussian Naive Bayes posteriors are overconfident on minority classes. One-vs-Rest isotonic regression maps them to calibrated probabilities:

\begin{equation}
\hat{P}(c_k\mid\mathbf{Z}) = f_{\text{iso}}^{(k)}\!\left(P_{\text{raw}}(c_k\mid\mathbf{Z})\right)
\end{equation}

where $f_{\text{iso}}^{(k)}$ is a monotone step function minimizing squared error on held-out calibration data. The Expected Calibration Error:

\begin{equation}
\text{ECE} = \sum_{m=1}^M \frac{|B_m|}{N}\,\bigl|\,\mathrm{acc}(B_m) - \mathrm{conf}(B_m)\,\bigr|
\end{equation}

\begin{tcolorbox}[resultbox]
\textbf{Phase 1 Calibration Result:} ECE = \textbf{0.038} (CI95: 0.020–0.057) against a target of $\le$0.05. Raw classifier ECE $\approx$ 0.12. Isotonic regression was selected over Platt scaling because the 8-class imbalanced distribution violates the sigmoid monotonicity assumption on raw Naive Bayes scores.
\end{tcolorbox}

\subsection{Pillar 5 — Expected-Loss Policy \& Active Evidence Acquisition}

The optimal operator action minimizes expected cost under the calibrated posterior:

\begin{equation}
a^* = \arg\min_{a \in \mathcal{A}} \sum_{k=1}^K \hat{P}(c_k\mid\mathbf{Z})\cdot\mathcal{L}(a,\,c_k)
\end{equation}

where $\mathcal{A} = \{\texttt{alert},\,\texttt{monitor},\,\texttt{ignore},\,\texttt{escalate\_human}\}$ and $\mathcal{L}$ is an asymmetric loss matrix penalizing cyber-attack false negatives heavily.

When $H(\hat{\mathbf{P}})$ is high between P2 and A1, the system requests the inverter sensor (Active Evidence Acquisition), maximizing mutual information gain:

\begin{equation}
\Delta H(S_{\text{inv}}) = H(\hat{\mathbf{P}}) - \mathbb{E}_y\!\left[H\!\left(\hat{\mathbf{P}}\mid S_{\text{inv}}=y\right)\right]
\end{equation}

Figures~\ref{fig:sensor_sensitivity} and~\ref{fig:active_evidence} illustrate the sensor selection rationale.

\begin{figure}[H]
  \centering
  \begin{subfigure}[b]{0.48\linewidth}
    \includegraphics[width=\linewidth]{figures/fig05_sensor_sensitivity.png}
    \caption{Voltage sensor sensitivity at PCC under varying noise and line impedance. The shaded region marks the nominal operating range; the detection margin narrows with higher measurement noise.}
    \label{fig:sensor_sensitivity}
  \end{subfigure}
  \hfill
  \begin{subfigure}[b]{0.48\linewidth}
    \includegraphics[width=\linewidth]{figures/fig06_active_evidence.png}
    \caption{Active Evidence Acquisition: entropy reduction when querying the hidden inverter sensor. Resolves P2 vs.\ A1 ambiguity that cannot be settled from PCC measurements alone.}
    \label{fig:active_evidence}
  \end{subfigure}
  \caption{Goal 2 — Uncertainty quantification and active measurement selection.}
\end{figure}

\subsection{Pillar 6 — Governance \& Audit Chain (Goal 3)}

Every alert produces a tamper-evident log entry in SQLite:
\begin{itemize}
  \item \textbf{Content:} Timestamp, detected class, calibrated confidence, action taken, decision rationale, SHAP attributions, $\mathbf{Z}$ residuals
  \item \textbf{Chain integrity:} SHA-256 hash of the current record includes the hash of the previous record — any modification breaks the chain
  \item \textbf{Human review rate:} 72.1\% of alerts are escalated to the human operator
\end{itemize}

\subsection{Pillar 7 — Noosphere Collective Defense (Goal 4)}

Noosphere implements peer-to-peer federated knowledge sharing across a \textbf{multi-node network} of energy systems — not limited to a single partner. The protocol is:

\begin{enumerate}
  \item \textbf{Compute:} Class-conditional feature statistics $(\boldsymbol{\mu}_k, \boldsymbol{\sigma}_k, n_k)$ per fault class
  \item \textbf{Protect:} Add Laplace noise calibrated to $\varepsilon=1.0$ differential privacy; suppress signatures with $n_k < 5$ (k-anonymity)
  \item \textbf{Export:} Pydantic-validated \texttt{IncidentSignatureV2} packets (JSON-serializable, schema v2.0)
  \item \textbf{Share:} Push to peer nodes (factory, hospital, smart building, rural microgrid, regional grid operator)
  \item \textbf{Import:} Receiving nodes check PSI + Wasserstein domain shift, then verify sensor compatibility
  \item \textbf{Merge:} Compatible signatures update the receiving node's Bayesian prior with weight $w=0.3$
\end{enumerate}

\begin{equation}
\tilde{\boldsymbol{\mu}}_k = \boldsymbol{\mu}_k + \mathrm{Lap}\!\left(\frac{\Delta f}{\varepsilon}\right),\quad
\Delta f = \frac{\max\|\mathbf{Z}\|_1}{n_k}
\end{equation}

Figure~\ref{fig:noosphere} shows the multi-node transfer results for the campus-to-factory link (Phase 1 demonstration).

\begin{figure}[H]
  \centering
  \includegraphics[width=0.95\linewidth]{figures/fig08_noosphere_transfer.png}
  \caption{Noosphere federated knowledge transfer. \textbf{Phase 1 result (campus $\to$ factory node):} 7 signatures exported; 4 imported with $\varepsilon=1.0$ DP; 3 rejected because the factory node lacks the \texttt{inv} and \texttt{vinv} sensor channels. \textbf{Phase 2 target:} Extend to a 5-node peer network (campus, hospital, industrial park, smart building, rural microgrid) with bidirectional exchange and a federated global threat registry.}
  \label{fig:noosphere}
\end{figure}

% ─────────────────────────────────────────────────────────────────────────────
\section{Experimental Results}
% ─────────────────────────────────────────────────────────────────────────────

\subsection{Experimental Protocol}

All experiments use 5 deterministic seeds $S \in \{22, 33, 44, 55, 66\}$. Each seed injects 96 test events across 8 fault classes. Calibration uses 80 independent events. Test and calibration sets are strictly disjoint.

\subsection{Goal 1 — Detection Benchmark}

\begin{figure}[H]
  \centering
  \includegraphics[width=0.98\linewidth]{figures/fig04_detection_benchmark.png}
  \caption{Detection benchmark across five methods (5 seeds, 95\% CI). \textbf{Left:} Event recall by method. Raw IsoForest fails completely (3.5\%) because seasonal cycles dominate unnormalized telemetry. Physics-only reaches 49.0\%. CUSUM achieves high raw recall (83.9\%) but at an unacceptable 57.8\% FPR ($\approx$184 false alarms/day). EnergyMind fusion achieves the highest balanced recall (58.1\%) with 91.2\% precision and only 0.173 FA/day. \textbf{Center:} Precision-Recall scatter. \textbf{Right:} Per-class event recall — grid events and sensor faults reach 100\%; cyber-attack recall (45\%) is limited by the hidden inverter sensor.}
  \label{fig:benchmark}
\end{figure}

\begin{table}[H]
\centering
\caption{Comprehensive Benchmark Results (5-Seed Mean with 95\% CI)}
\begin{tabularx}{\linewidth}{Xcccc}
\toprule
\textbf{Method} & \textbf{Event Recall} & \textbf{Precision} & \textbf{FPR} & \textbf{FA / Day}\\
\midrule
A: IsoForest (Raw) & 0.035 [−0.003, 0.074] & 0.463 & 0.0006 & $\approx$2.8\\
B: Physics Only & 0.490 [0.477, 0.503] & 0.922 & 0.0026 & 0.62\\
C: IsoForest + Residuals & 0.450 [0.415, 0.485] & 0.908 & 0.0022 & 0.53\\
D: CUSUM (Standalone) & 0.840 [0.810, 0.869] & 0.121 & 0.5777 & 184.2\\
\midrule
\textbf{EnergyMind (Fusion)} & \textbf{0.581 [0.545, 0.617]} & \textbf{0.912} & \textbf{0.0037} & \textbf{0.173}\\
\bottomrule
\end{tabularx}
\end{table}

\subsection{Goal 2 — Calibration \& Selective Prediction}

\begin{figure}[H]
  \centering
  \includegraphics[width=0.72\linewidth]{figures/fig07_risk_coverage.png}
  \caption{Selective prediction risk-coverage curve. At the operating point (49.4\% coverage), EnergyMind achieves 0\% error rate on decided cases — all 5 seeds confirm 100\% diagnostic accuracy. The tradeoff: 50.6\% of alerts are flagged as uncertain and routed to the operator. When abstention is disabled (100\% coverage), accuracy drops to 86.1\%, demonstrating the honest value of selective prediction.}
  \label{fig:risk_coverage}
\end{figure}

\begin{table}[H]
\centering
\caption{Goal 2 — Uncertainty Quantification Metrics}
\begin{tabular}{lll}
\toprule
\textbf{Metric} & \textbf{Value} & \textbf{Target / Interpretation}\\
\midrule
ECE (calibrated) & 0.038 (CI95: 0.020–0.057) & $\le$0.05 ✓\\
ECE (raw, uncalibrated) & $\approx$0.12 & Baseline before calibration\\
Brier score & 0.036 (CI95: 0.033–0.039) & Lower is better\\
Abstention rate & 50.6\% & Refuses when uncertain\\
Accuracy (decided cases) & 100.0\% (5/5 seeds) & No wrong diagnoses when system commits\\
Accuracy (forced, no abstention) & 86.1\% & Accuracy without selective prediction\\
Human review rate & 72.1\% & Proportion escalated to operator\\
\bottomrule
\end{tabular}
\end{table}

\subsection{Goal 4 — Noosphere Knowledge Transfer}

\begin{table}[H]
\centering
\caption{Noosphere Transfer Results — Campus Node $\to$ Factory Node (Phase 1)}
\begin{tabular}{lll}
\toprule
\textbf{Fault Class} & \textbf{Status} & \textbf{Compatibility Check}\\
\midrule
\texttt{pv\_fault} & \textcolor{AccentGreen}{\checkmark Imported} & $w=0.3$, $\varepsilon=1.0$, sensors compatible\\
\texttt{grid\_event} & \textcolor{AccentRed}{$\times$ Rejected} & Missing sensor: \texttt{vinv}\\
\texttt{load\_peak} & \textcolor{AccentGreen}{\checkmark Imported} & $w=0.3$, $\varepsilon=1.0$, sensors compatible\\
\texttt{pv\_sensor\_fault} & \textcolor{AccentRed}{$\times$ Rejected} & Missing sensor: \texttt{inv}\\
\texttt{voltage\_sensor\_fault} & \textcolor{AccentGreen}{\checkmark Imported} & $w=0.3$, $\varepsilon=1.0$, sensors compatible\\
\texttt{bess\_sensor\_fault} & \textcolor{AccentGreen}{\checkmark Imported} & $w=0.3$, $\varepsilon=1.0$, sensors compatible\\
\texttt{coordinated\_falsification} & \textcolor{AccentRed}{$\times$ Rejected} & Missing sensor: \texttt{inv}\\
\bottomrule
\end{tabular}
\end{table}

\begin{tcolorbox}[warnbox]
\small\textbf{Important:} The Noosphere protocol is designed for \emph{multi-node} peer exchange, not just bilateral transfer. Phase 1 demonstrates the protocol on the campus-to-factory link as a proof of concept. Phase 2 will expand to a 5-node network: factory, hospital, smart commercial building, rural agri-solar microgrid, and regional grid operator.
\end{tcolorbox}

% ─────────────────────────────────────────────────────────────────────────────
\section{Honest Engineering — Limitations}
% ─────────────────────────────────────────────────────────────────────────────

\begin{table}[H]
\centering
\caption{Known Limitations and Phase 2 Mitigations}
\begin{tabularx}{\linewidth}{lXX}
\toprule
\textbf{Limitation} & \textbf{Impact} & \textbf{Phase 2 Mitigation}\\
\midrule
\texttt{load\_peak} recall = 0\% & AR(1) model absorbs short peaks into noise & Dedicated load forecasting (LSTM or Gaussian Process)\\
Cyber attack recall = 45\% & Requires hidden \texttt{inv} sensor & Real inverter telemetry from campus SCADA\\
Abstention rate 50.6\% & High for operational use & More calibration events ($N>300$)\\
Human review rate 72.1\% & Conservative in Phase 1 & Reduce to $\le$40\% target with more data\\
CUSUM FPR = 57.7\% & Cannot use standalone & Used as additive channel only (never primary)\\
Simulated BESS telemetry & No real BESS measurements & Hardware-in-the-loop integration (Phase 2)\\
1-node Noosphere demo & Only campus-factory tested & 5-node peer network (Phase 2)\\
\bottomrule
\end{tabularx}
\end{table}

% ─────────────────────────────────────────────────────────────────────────────
\section{Phase 2 Roadmap}
% ─────────────────────────────────────────────────────────────────────────────

\begin{table}[H]
\centering
\caption{Phase 2 Engineering Milestones}
\begin{tabularx}{\linewidth}{llX}
\toprule
\textbf{Priority} & \textbf{Milestone} & \textbf{Target Metric}\\
\midrule
High & Real campus SCADA streaming telemetry & Inference latency $< 250$\,ms\\
High & Expand calibration data ($N > 300$ per class) & ECE $< 0.03$, abstention $< 30\%$\\
High & 5-node Noosphere peer network & Bidirectional exchange, federated global registry\\
Medium & LightGBM / XGBoost multi-class diagnoser & Forced accuracy $> 92\%$\\
Medium & Conformal prediction intervals on regression & 95\% empirical coverage guarantee\\
Low & IEC 61850 GOOSE/SV parsing & Grid standards compliance\\
Low & Hardware-in-the-loop (OPAL-RT) testing & Real-time grid validation\\
\bottomrule
\end{tabularx}
\end{table}

% ─────────────────────────────────────────────────────────────────────────────
\section{Repository \& Deliverables}
% ─────────────────────────────────────────────────────────────────────────────

\begin{table}[H]
\centering
\caption{Phase 1 Required Deliverables — Checklist}
\begin{tabular}{lll}
\toprule
\textbf{Deliverable} & \textbf{File / Location} & \textbf{Status}\\
\midrule
Project Proposal (this report) & \texttt{docs/REPORT\_PHASE1.tex} & ✓ Complete\\
Scope definition & Section 2 of this report & ✓ Complete\\
Objectives & Section 3 of this report & ✓ Complete\\
Data / anomaly simulation plan & Section 4 of this report & ✓ Complete\\
Preliminary design for each Goal & Section 5 of this report & ✓ Complete\\
GitHub Repository & \url{https://github.com/anonymous-tsyp14/energymind} & ✓ Public\\
Dataset & \texttt{data/raw/} + \texttt{data/reference/} & ✓ Committed\\
Code developed in Phase 1 & \texttt{src/energymind/} + \texttt{scripts/} & ✓ Committed\\
Unit tests & \texttt{tests/test\_core.py} & ✓ 21/21 PASS\\
Monitoring platform & \texttt{app/dashboard.py} (FastAPI) & ✓ Running\\
Publication figures & \texttt{outputs/fig01--fig08.png} & ✓ Generated\\
\bottomrule
\end{tabular}
\end{table}

\textbf{To reproduce all results:}
\begin{verbatim}
python -m venv .venv && .venv\Scripts\activate
pip install -r requirements.txt
python scripts/run_v2.py          # 5 seeds, ~5 min
python scripts/make_figures_v2.py # 8 publication figures
pytest tests/ -v                  # 21/21 tests
uvicorn app.dashboard:app --port 8000  # Live SCADA platform
\end{verbatim}

% ─────────────────────────────────────────────────────────────────────────────
%  REFERENCES
% ─────────────────────────────────────────────────────────────────────────────

\begin{thebibliography}{16}

\bibitem{ieee1547}
IEEE Std 1547-2018, \emph{Standard for Interconnection and Interoperability of Distributed Energy Resources with Associated Electric Power Systems Interfaces}, IEEE, 2018. \url{https://doi.org/10.1109/IEEESTD.2018.8332112}

\bibitem{pandapower}
L. Thurner, A. Scheidler, F. Schäfer \emph{et al.}, ``pandapower: An Open-Source Python Tool for Convenient Modeling, Analysis, and Optimization of Electric Power Systems,'' \emph{IEEE Trans. Power Systems}, vol. 33, no. 6, pp. 6510--6521, 2018. \url{https://doi.org/10.1109/TPWRS.2018.2829021}

\bibitem{isoforest}
F. T. Liu, K. M. Ting, Z.-H. Zhou, ``Isolation Forest,'' \emph{Proc. 8th IEEE ICDM}, pp. 413--422, 2008. \url{https://doi.org/10.1109/ICDM.2008.17}

\bibitem{shap}
S. M. Lundberg, S.-I. Lee, ``A Unified Approach to Interpreting Model Predictions,'' \emph{NeurIPS}, vol. 30, 2017. \url{https://arxiv.org/abs/1705.07874}

\bibitem{zadrozny}
B. Zadrozny, C. Elkan, ``Transforming Classifier Scores into Accurate Multiclass Probability Estimates,'' \emph{Proc. 8th ACM SIGKDD}, pp. 694--699, 2002. \url{https://doi.org/10.1145/775047.775151}

\bibitem{dwork_dp}
C. Dwork, ``Differential Privacy,'' \emph{Proc. ICALP}, LNCS 4052, pp. 1--12, 2006. \url{https://doi.org/10.1007/11787006_1}

\bibitem{selective_pred}
Y. Geifman, R. El-Yaniv, ``Selective Classification for Deep Neural Networks,'' \emph{NeurIPS}, 2017. \url{https://arxiv.org/abs/1705.08500}

\bibitem{pvgis}
T. Huld \emph{et al.}, ``A New Solar Radiation Database for Estimating PV Performance in Europe and Africa,'' \emph{Solar Energy}, vol. 86, pp. 1803--1815, 2012. \url{https://re.jrc.ec.europa.eu/pvg_tools/en/}

\bibitem{nasa_power}
P. W. Stackhouse \emph{et al.}, ``NASA POWER Project,'' NASA Langley Research Center, 2023. \url{https://power.larc.nasa.gov/}

\bibitem{simbench}
S. Meinecke \emph{et al.}, ``SimBench: A Benchmark Dataset of Electric Power Systems,'' \emph{Energies}, vol. 13, no. 12, 2020. \url{https://doi.org/10.3390/en13123412}

\bibitem{doe_bldg}
N. Wang \emph{et al.}, ``Commercial Reference Building Models of the National Building Stock,'' NREL/TP-5500-46861, 2011. \url{https://www.energy.gov/eere/buildings/commercial-reference-buildings}

\bibitem{uah_ics}
T. H. Morris, W. Gao, ``Industrial Control System Cyber Attacks,'' \emph{Proc. ICS-CSR}, 2014. \url{https://sites.google.com/a/uah.edu/tommy-morris-uah/ics-data-sets}

\bibitem{lightgbm}
G. Ke \emph{et al.}, ``LightGBM: A Highly Efficient Gradient Boosting Decision Tree,'' \emph{NeurIPS}, vol. 30, pp. 3146--3154, 2017.

\bibitem{niculescu}
A. Niculescu-Mizil, R. Caruana, ``Predicting Good Probabilities with Supervised Learning,'' \emph{Proc. ICML}, pp. 625--632, 2005.

\bibitem{cusum_basseville}
M. Basseville, I. V. Nikiforov, \emph{Detection of Abrupt Changes: Theory and Application}, Prentice Hall, 1993.

\bibitem{geifman_elyaniv}
Y. Geifman, R. El-Yaniv, ``Selective Prediction-Set Models with Coverage Guarantees,'' \emph{NeurIPS}, 2019. \url{https://arxiv.org/abs/1901.08903}

\end{thebibliography}

\end{document}
